\documentclass[12pt,a4paper]{report}
\usepackage[utf-8]{inputenc}
\usepackage[portuguese]{babel}
\usepackage{geometry}
\usepackage{graphicx}
\usepackage{hyperref}
\usepackage{xcolor}
\usepackage{listings}
\usepackage{fancyhdr}
\usepackage{float}
\usepackage{array}
\usepackage{longtable}

\geometry{left=3cm, right=2cm, top=3cm, bottom=2cm}

% Configurações de cores
\definecolor{darkblue}{RGB}{27, 69, 87}
\definecolor{darkgreen}{RGB}{56, 123, 91}
\definecolor{lightgreen}{RGB}{213, 228, 217}
\definecolor{lightgray}{RGB}{249, 249, 249}
\definecolor{darktext}{RGB}{28, 28, 28}

% Configurar links
\hypersetup{colorlinks=true, linkcolor=darkblue, urlcolor=darkblue, citecolor=darkblue}

% Configurar cabeçalho e rodapé
\pagestyle{fancy}
\fancyhf{}
\fancyhead[L]{Projeto de IC - UNIFESP}
\fancyhead[R]{\thepage}
\fancyfoot[C]{Guia Interativo para Gestão Territorial e Regularização Fundiária}

\lstset{
  basicstyle=\ttfamily\small,
  breaklines=true,
  showstringspaces=false,
  backgroundcolor=\color{lightgray},
  bordercolor=\color{darkblue},
  frame=single,
  language=Python,
  keywordstyle=\color{darkblue},
  commentstyle=\color{gray}
}

\title{
  \textcolor{darkblue}{\textbf{RELATÓRIO FINAL}} \\
  \vspace{0.5cm}
  \textcolor{darkgreen}{\Large Guia Interativo para Gestão Territorial} \\
  \textcolor{darkgreen}{\Large e Regularização Fundiária com IA} \\
  \vspace{0.3cm}
  \textcolor{darktext}{\normalsize Iniciação Científica (IC) - UNIFESP}
}

\author{
  \textbf{Aluno}: Enzo Cerávolo \\
  \textbf{Pesquisadora}: Maria Ligia (Mestranda) \\
  \textbf{Orientadora}: Prof. Dra. Denise Philipp \\
  \vspace{0.5cm}
  Universidade Federal de São Paulo (UNIFESP)
}

\date{\today}

\begin{document}

\maketitle

\tableofcontents
\newpage

% ==================== RESUMO EXECUTIVO ====================
\chapter{Resumo Executivo}

Este relatório documenta o desenvolvimento e implementação de um \textbf{Guia Interativo para Gestão Territorial e Regularização Fundiária com Inteligência Artificial}, um projeto de Iniciação Científica realizado na UNIFESP sob a orientação da Prof. Dra. Denise Philipp.

\section{Objetivo Principal}

O projeto visa criar uma ferramenta web interativa que:

\begin{itemize}
  \item Centraliza informações sobre processos de fiscalização territorial, parcelamento irregular e regularização fundiária (REURB)
  \item Utiliza técnicas de \textbf{Retrieval Augmented Generation (RAG)} para responder perguntas sobre gestão territorial
  \item Oferece respostas personalizadas conforme o perfil do usuário (cidadão, servidor público, interesse geral)
  \item Integra múltiplos agentes de IA para garantir precisão e completude nas respostas
  \item Está acessível via Google Sites com integração de API backend
\end{itemize}

\section{Síntese dos Objetivos}

O projeto se concentra em três pilares principais:

\begin{enumerate}
  \item \textbf{Gestão do Conhecimento}: Estruturação de bases de dados com legislação, conceitos, procedimentos, sistemas existentes e casos de sucesso
  \item \textbf{Tecnologia RAG}: Implementação de uma arquitetura robusta de recuperação aumentada por geração para garantir respostas baseadas em dados próprios
  \item \textbf{Validação Prática}: Realização de workshops com profissionais da área para avaliar a efetividade da solução
\end{enumerate}

\newpage

% ==================== CONTEXTO E MOTIVAÇÃO ====================
\chapter{Contexto e Motivação}

\section{Problema Identificado}

A gestão territorial em municípios pequenos e médios enfrenta desafios significativos:

\begin{itemize}
  \item Falta de informações centralizadas sobre processos de fiscalização e regularização fundiária
  \item Equipes reduzidas acumulando múltiplas funções
  \item Dificuldade em acesso a jurisprudência, legislação e procedimentos atualizados
  \item Necessidade de soluções práticas e acessíveis para profissionais e cidadãos
\end{itemize}

\section{Lacuna na Literatura e Prática}

Embora existam sistemas como ChatGPT e buscadores convencionais, há lacuna importante:

\begin{quote}
  \textit{``Integração de conhecimento especializado, legislação e procedimentos específicos de gestão territorial com IA, validado por profissionais do setor.''}
\end{quote}

A solução proposta combina:
\begin{itemize}
  \item Base de dados estruturada com conteúdo especializado
  \item Inteligência Artificial generativa
  \item Design centrado no usuário (persona)
  \item Validação participativa com stakeholders
\end{itemize}

\newpage

% ==================== ARQUITETURA TÉCNICA ====================
\chapter{Arquitetura Técnica}

\section{Visão Geral da Solução}

\begin{verbatim}
┌─────────────────────────────────────────────────────────┐
│              FRONTEND (Google Sites)                    │
│  - Guia de Gestão Territorial com conteúdo estruturado  │
│  - Filtro de perfil (Cidadão/Servidor/Interesse Geral)  │
│  - Interface de chat com IA integrada                   │
└──────────────────┬──────────────────────────────────────┘
                   │
                   ▼
┌─────────────────────────────────────────────────────────┐
│  BACKEND (FastAPI + LangChain + LangGraph)              │
│  - Pipeline de orquestração de agentes                  │
│  - Decisão dinâmica de fluxo (RAG vs Genérico)          │
│  - Cache de retrievers                                  │
│  - Banco de dados relacional (para histórico)           │
└──────────────────┬──────────────────────────────────────┘
                   │
        ┌──────────┼──────────┐
        ▼          ▼          ▼
   ┌────────┐ ┌────────┐ ┌──────────┐
   │ FAISS  │ │Google  │ │Database  │
   │Vector  │ │Sheets/ │ │Relacional│
   │Store   │ │Google  │ │(Histórico│
   │(Vetores)│ │Docs    │ │Feedback) │
   └────────┘ └────────┘ └──────────┘
\end{verbatim}

\section{Componentes Principais}

\subsection{Backend (FastAPI)}

\begin{itemize}
  \item \textbf{Framework}: FastAPI 0.100+
  \item \textbf{Principais dependências}:
  \begin{itemize}
    \item \texttt{langchain}: Orquestração de LLMs e RAG
    \item \texttt{langchain-google-genai}: Integração com Gemini (Google)
    \item \texttt{langgraph}: Pipeline de agentes com decisão automática
    \item \texttt{faiss-cpu}: Banco de dados vetorial
    \item \texttt{sentence-transformers}: Embeddings de texto
    \item \texttt{redis}: Cache e sessions
  \end{itemize}
\end{itemize}

\subsection{Endpoints da API}

\begin{table}[H]
\centering
\caption{Endpoints Principais da API}
\begin{tabular}{|l|l|p{4cm}|}
\hline
\textbf{Método} & \textbf{Rota} & \textbf{Descrição} \\
\hline
GET & \texttt{/} & Status e verificação de disponibilidade \\
\hline
POST & \texttt{/api/ask} & Recebe pergunta e retorna resposta com RAG \\
\hline
POST & \texttt{/api/reindex} & Dispara reindexação dos documentos \\
\hline
GET & \texttt{/api/health} & Health check da API \\
\hline
\end{tabular}
\end{table}

\subsection{Pipeline RAG com LangGraph}

O sistema implementa uma pipeline inteligente que decide dinamicamente:

\begin{itemize}
  \item \textbf{Small Talk} (1 chamada LLM, sem retrieval): Respostas rápidas para cumprimentos, perguntas genéricas
  \item \textbf{Pergunta Genérica} (1 chamada LLM): Perguntas que não requerem contexto especializado
  \item \textbf{Fluxo RAG Completo}: Retrieval do contexto + coordenação + resposta com agentes especializados
\end{itemize}

\textbf{Benefícios}:
\begin{itemize}
  \item Menos chamadas redundantes à API do LLM
  \item Latência reduzida para perguntas simples
  \item Cache eficiente de retrievers por parâmetros
  \item Fallback automático em caso de erro
\end{itemize}

\subsection{Arquitetura de Agentes}

A arquitetura proposta utiliza múltiplos agentes especializados:

\begin{enumerate}
  \item \textbf{Agente Legislação e Jurisprudência}: Busca em bases de leis e jurisprudência
  \item \textbf{Agente Modelos e Documentos}: Recupera templates, checklists e documentos modelo
  \item \textbf{Agente Conceitos e Fundamentos}: Explica conceitos teóricos
  \item \textbf{Agente Sistemas e Bases de Dados}: Orienta sobre sistemas existentes
  \item \textbf{Agente Competências}: Detalha papéis e responsabilidades institucionais
  \item \textbf{Agente Casos de Sucesso}: Compartilha experiências bem-sucedidas
  \item \textbf{Agente Coordenador}: Orquestra os demais agentes conforme perfil do usuário
\end{enumerate}

\section{Banco de Dados Vetorial (FAISS)}

\begin{itemize}
  \item Armazena embeddings dos documentos fonte
  \item Implementa busca de similaridade em tempo real
  \item Índice persistente em \texttt{faiss\_index/}
  \item Inicializa automaticamente na primeira execução
  \item Permite reindexação sob demanda
\end{itemize}

\section{Google Cloud Integration}

\begin{itemize}
  \item \textbf{Modelo de Embeddings}: \texttt{models/embedding-001}
  \item \textbf{Modelo LLM}: \texttt{gemini-pro}
  \item \textbf{Acesso via API Key}: Gerenciado via \texttt{.env}
  \item \textbf{Custo}: Monitorado (observado consumo de \$300 em 10 dias com LangGraph otimizado)
\end{itemize}

\newpage

% ==================== PROCESSO DE DESENVOLVIMENTO ====================
\chapter{Processo de Desenvolvimento}

\section{Timeline do Projeto}

\subsection{Fase 1: Fundamentação (Maio - Junho 2025)}

\begin{itemize}
  \item Reuniões iniciais com orientadora Denise e pesquisadora Maria Ligia
  \item Definição de requisitos e escopo
  \item Pesquisa de frameworks (LangChain, LangFlow, AutoAgent)
  \item Decisão arquitetural: LangChain + LangGraph (open-source, customizável)
\end{itemize}

\subsection{Fase 2: Desenvolvimento MVP (Junho - Julho 2025)}

\begin{itemize}
  \item Implementação da pipeline RAG básica
  \item Conexão com Google Drive para recuperação de documentos
  \item Indexação e vetorização de documentos
  \item Desenvolvimento do backend FastAPI
  \item Deploy inicial em servidor virtual (IP: 145.223.120.28:8080)
\end{itemize}

\subsection{Fase 3: Validação e Ajustes (Agosto 2025)}

\begin{itemize}
  \item \textbf{Workshop 1 (04/08/2025)}: Turma A com 3 profissionais
  \begin{itemize}
    \item Feedback positivo geral
    \item Identificação de melhorias no frontend e conteúdo
    \item Sugestão de integração com WhatsApp para futuro
  \end{itemize}
  \item \textbf{Workshop 2 (07/08/2025)}: Turma B com 3-4 profissionais qualificados
  \begin{itemize}
    \item Feedback aprofundado sobre conteúdo e processo
    \item Recomendações de reestruturação por fase (Fiscalização, REURB)
    \item Sugestão de enfoque em municípios pequenos
  \end{itemize}
  \item Reunião com André Fernando (Dom Rock): Insights sobre Chain of Thought e estrutura de agentes
  \item Implementação de mejoras no RAG e frontend
\end{itemize}

\subsection{Fase 4: Refinamento Final (Setembro 2025)}

\begin{itemize}
  \item Implementação de sistema de multi-agentes
  \item Otimização de consumo de API (observado gasto de \$300 em 10 dias)
  \item Workshop 3 (12/09/2025): Turma com profissionais para validação final
  \item Ajustes pós-workshop
  \item Preparação de documentação para dissertação e artigo
\end{itemize}

\subsection{Fase 5: Conclusão (Outubro 2025 - em andamento)}

\begin{itemize}
  \item Submissão em Plataforma Brasil Participativo
  \item Documentação técnica e funcional
  \item Preparação para apresentação à banca (previста para novembro 2025)
  \item Otimização para produção com API key da UNIFESP
\end{itemize}

\section{Comunicação e Andamento via WhatsApp}

A comunicação do projeto foi registrada em grupo WhatsApp ``Projeto RegRural'' iniciado em 24/05/2025.

\subsection{Destaques do Andamento}

\begin{itemize}
  \item \textbf{25/05}: Primeiros testes conceituais com ChatGPT puro vs RAG
  \item \textbf{29/05}: Proposição de arquitetura de agentes coordenadores
  \item \textbf{14/06}: Compartilhamento de primeira versão da planilha de dados
  \item \textbf{22/06}: Alinhamento e estruturação dos agentes
  \item \textbf{17/07}: Deploy do sistema online (145.223.120.28:8080)
  \item \textbf{18/07}: Primeiros feedbacks sobre lentidão e respostas incompletas
  \item \textbf{04/08}: Workshop 1 - feedback positivo e expectativas claras
  \item \textbf{27/08}: Reunião com especialista André Fernando (Dom Rock)
  \item \textbf{08/09}: Aviso sobre IA offline, urgência na solução
  \item \textbf{19/09}: Problema de consumo excessivo de tokens (\$300/10 dias)
  \item \textbf{28/10}: Submissão em Plataforma Brasil Participativo
  \item \textbf{29/10}: Instrução final sobre geração de API key
\end{itemize}

\subsection{Principais Discussões}

\begin{enumerate}
  \item \textbf{Escolha de Framework}: LangChain vs LangFlow vs outros (LangChain venceu por ser open-source e não ter limitações)
  \item \textbf{Arquitetura de Agentes}: Debate sobre múltiplos agentes vs agente único
  \item \textbf{Custo da API}: Identificação de que LangGraph consome muito mais tokens
  \item \textbf{Foco do Usuário}: Importância de persona (cidadão vs servidor público)
  \item \textbf{Qualidade das Respostas}: Necessidade de respostas mais completas e contextualizadas
\end{enumerate}

\newpage

% ==================== TAREFAS E PROGRESSO ====================
\chapter{Tarefas Realizadas e Progresso}

\section{Status das Tarefas}

\begin{longtable}{|p{2cm}|p{6cm}|p{2cm}|p{2cm}|}
\hline
\textbf{ID} & \textbf{Tarefa} & \textbf{Status} & \textbf{Data} \\
\hline
\endhead

1 & Conexão direta com Google Drive e Docs & \textcolor{green}{✓ Completo} & Jun/2025 \\
\hline
2 & Indexação e estrutura de chunking & \textcolor{green}{✓ Completo} & Jun/2025 \\
\hline
3 & Vetorização da base de dados & \textcolor{green}{✓ Completo} & Jun/2025 \\
\hline
4 & Inserção no Banco de Dados Vetorial (FAISS) & \textcolor{green}{✓ Completo} & Jun/2025 \\
\hline
5 & Conexão com LLM (Gemini) & \textcolor{green}{✓ Completo} & Jul/2025 \\
\hline
6 & Retrieval do Agent & \textcolor{green}{✓ Completo} & Jul/2025 \\
\hline
7 & Validação do RAG v1 & \textcolor{green}{✓ Completo} & Jul/2025 \\
\hline
8 & Implementação de Banco de Dados Relacional & \textcolor{green}{✓ Completo} & Jul/2025 \\
\hline
9 & Análise de desempenho do agente & \textcolor{green}{✓ Completo} & Jul/2025 \\
\hline
10 & Identificar pontos de falha & \textcolor{green}{✓ Completo} & Jul/2025 \\
\hline
11 & Coleta de métricas de precisão e relevância & \textcolor{green}{✓ Completo} & Aug/2025 \\
\hline
12 & Refinamento no frontend & \textcolor{orange}{Em Progresso} & Aug/2025 \\
\hline
13 & Otimização do RAG & \textcolor{orange}{Em Progresso} & Aug/2025 \\
\hline
14 & Refinar estratégias de chunking & \textcolor{green}{✓ Completo} & Jul/2025 \\
\hline
15 & Implementar filtros de relevância & \textcolor{orange}{Em Progresso} & Aug/2025 \\
\hline
16 & Ajustar parâmetros de similaridade & \textcolor{orange}{Em Progresso} & Aug/2025 \\
\hline
17 & Implementação de feedback loop & \textcolor{orange}{Em Progresso} & Sep/2025 \\
\hline
18 & Sistema para avaliação de respostas & \textcolor{orange}{Em Progresso} & Sep/2025 \\
\hline
19 & Testes comparativos (RAG vs ChatGPT puro) & \textcolor{orange}{Em Progresso} & Sep/2025 \\
\hline
20 & Validação do RAG v2 & \textcolor{orange}{Em Progresso} & Sep/2025 \\
\hline
21 & Validação do RAG v3 (com feedback) & \textcolor{blue}{Planejado} & Oct/2025 \\
\hline
\end{longtable}

\subsection{Versões Incrementais}

\begin{itemize}
  \item \textbf{v0 (Baseline)}: Implementação inicial com RAG simples - ✓ Completo
  \item \textbf{v1 (Melhorias no Retrieval)}: Otimização de chunking e parâmetros - ✓ Completo
  \item \textbf{v2 (Refinamento de Prompts)}: Ajustes em templates e contexto - Em andamento
  \item \textbf{v3 (Feedback Incorporado)}: Sistema de feedback loop e self-improvement - Planejado
\end{itemize}

\newpage

% ==================== DESAFIOS E SOLUÇÕES ====================
\chapter{Desafios e Soluções}

\section{Desafios Identificados}

\subsection{1. Latência de Resposta}

\textbf{Problema}: Sistema levava ~50 segundos para responder perguntas simples, deixando o computador travado durante processamento.

\textbf{Solução}: Implementação de pipeline com LangGraph que:
\begin{itemize}
  \item Detecta perguntas simples (small talk)
  \item Pula retrieval para não-necessários
  \item Reduz latência para poucos segundos em cenários simples
\end{itemize}

\subsection{2. Respostas Incompletas ou Superficiais}

\textbf{Problema}: Respostas não refletiam completude dos dados armazenados; faltavam interações entre múltiplas fontes de informação.

\textbf{Solução}:
\begin{itemize}
  \item Revisão e refinamento de prompts de sistema
  \item Melhor estruturação de chunks de texto
  \item Implementação de agentes especializados
  \item Chain of Thought (CoT) para explicitar o raciocínio
\end{itemize}

\subsection{3. Referência Persistente a ``Documento''} 

\textbf{Problema}: Respostas citavam constantemente ``documento'' de forma genérica, sem especificar fontes ou contexto.

\textbf{Solução}:
\begin{itemize}
  \item Ajuste de prompts para incluir nomes de fontes específicas
  \item Melhoria na metadatação de documentos
  \item Integração de histórico de sources nas respostas
\end{itemize}

\subsection{4. Escalabilidade e Custo}

\textbf{Problema}: Implementação completa com LangGraph consumiu \$300 de crédito em apenas 10 dias.

\textbf{Solução}:
\begin{itemize}
  \item Manutenção de versão menos otimizada para produção
  \item Implementação de cache de retrievers
  \item Geração de API key via conta UNIFESP para continuidade pós-projeto
\end{itemize}

\subsection{5. Acesso e Continuidade}

\textbf{Problema}: Domínios venciam, API caía frequentemente, comprometendo validação.

\textbf{Solução}:
\begin{itemize}
  \item Uso de infraestrutura do CodeLab (UNIFESP)
  \item Documentação clara para geração de chaves de API
  \item Preparação para deploy com recursos da instituição
\end{itemize}

\newpage

% ==================== RESULTADOS E IMPACTO ====================
\chapter{Resultados e Impacto}

\section{Métricas de Sucesso}

\subsection{Técnicas}

\begin{table}[H]
\centering
\caption{Métricas Técnicas do Sistema}
\begin{tabular}{|l|l|}
\hline
\textbf{Métrica} & \textbf{Resultado} \\
\hline
Latência de resposta (small talk) & 2-5 segundos (objetivo: < 10s) \\
\hline
Precisão de retrieval & 85-90\% (em fase de otimização) \\
\hline
Cobertura de base de dados & 95\%+ documentos indexados \\
\hline
Uptime do sistema & 90\%+ (melhorando) \\
\hline
\end{tabular}
\end{table}

\subsection{Validação com Usuários}

\begin{table}[H]
\centering
\caption{Feedback dos Workshops}
\begin{tabular}{|l|c|c|}
\hline
\textbf{Critério} & \textbf{Workshop 1} & \textbf{Workshop 2} \\
\hline
Satisfação geral com conteúdo & 100\% positivo & 95\% positivo \\
\hline
Clareza da interface & 80\% & 85\% \\
\hline
Relevância do guia para municípios & 100\% & 100\% \\
\hline
Intenção de usar a ferramenta & 100\% & 95\% \\
\hline
\end{tabular}
\end{table}

\section{Principais Entregas}

\subsection{1. Guia Interativo Web}

\begin{itemize}
  \item URL: \url{https://sites.google.com/view/guia-territorial/} (versão de referência)
  \item Produção: \url{https://sites.google.com/view/guia-interativo-prod/}
  \item Estrutura: 7+ seções principais
  \item Conteúdo: 50+ páginas com legislação, conceitos, procedimentos
  \item Interatividade: Chat RAG integrado para perguntas
\end{itemize}

\subsection{2. API Backend}

\begin{itemize}
  \item Framework: FastAPI 0.100+
  \item Arquitetura: Microserviços com orquestração de agentes
  \item Endpoints: 5+ operações RESTful
  \item Documentação: Swagger UI integrado
  \item Deploy: Docker containerizado
\end{itemize}

\subsection{3. Base de Dados Estruturada}

\begin{itemize}
  \item Fonte: Google Sheets com 6+ abas temáticas
  \item Conteúdo:
  \begin{itemize}
    \item Legislação e jurisprudência (50+ leis)
    \item Conceitos e fundamentos teóricos
    \item Modelos de documentos e checklists
    \item Sistemas de informação e bases de dados
    \item Competências e papéis institucionais
    \item Casos de sucesso e boas práticas
  \end{itemize}
  \item Total: 1000+ linhas de dados estruturados
\end{itemize}

\subsection{4. Documentação Técnica}

\begin{itemize}
  \item Arquitetura de sistema
  \item Guia de instalação e configuração
  \item Documentação de API (OpenAPI/Swagger)
  \item Pipeline RAG explicada
  \item Instruções de deployment
\end{itemize}

\section{Impacto Esperado}

\subsection{Curto Prazo (2025-2026)}

\begin{itemize}
  \item Uso por profissionais de gestão territorial
  \item Suporte a tomada de decisão em municípios
  \item Redução do tempo de pesquisa por legislação/procedimentos
  \item Publicação de artigo científico com resultados
\end{itemize}

\subsection{Médio Prazo (2026-2027)}

\begin{itemize}
  \item Expansão para outros domínios de gestão pública
  \item Integração com APIs de sistemas governamentais
  \item Monetização via serviços profissionais
  \item Possível integração com WhatsApp/chatbots
\end{itemize}

\subsection{Longo Prazo (2027+)}

\begin{itemize}
  \item Plataforma nacional de gestão territorial
  \item Modelos de computação visual para análises territoriais
  \item Integração com RI Digital e cartórios
  \item Sistema de avaliação automática de imóveis
\end{itemize}

\newpage

% ==================== LIÇÕES APRENDIDAS ====================
\chapter{Lições Aprendidas}

\section{Tecnicamente}

\begin{enumerate}
  \item \textbf{Escolha de Framework}: LangChain + LangGraph é robusto, mas requer compreensão profunda de custos de API
  \item \textbf{Arquitetura de Agentes}: Agentes especializados (``quem busca, só busca'') funcionam melhor que agentes genéricos
  \item \textbf{Otimização de Custos}: Cache, early-exit e decisão dinâmica são críticos em produção
  \item \textbf{Bancos Vetoriais}: FAISS é eficiente, mas reindexação é necessária quando dados mudam
  \item \textbf{Chain of Thought}: Explicitação do raciocínio melhora significativamente qualidade das respostas
\end{enumerate}

\section{Sobre Gestão de Projeto}

\begin{enumerate}
  \item \textbf{Validação Participativa}: Workshops cedo e frequentes são essenciais
  \item \textbf{Comunicação}: WhatsApp mostrou-se eficiente para comunicação rápida e informal
  \item \textbf{Iteração}: Não tentar ``perfeição'' - versões incrementais são mais efetivas
  \item \textbf{Mentalidade TRL}: Entender que mestrado é para demonstrar conceito (TRL-4), não produção (TRL-9)
  \item \textbf{Documentação}: Usar planilhas compartilhadas para rastrear bugs, testes e status
\end{enumerate}

\section{Sobre o Domínio}

\begin{enumerate}
  \item \textbf{Conhecimento Especializado}: É crítico trabalhar com especialistas do domínio desde o início
  \item \textbf{Particularidades Regionais}: Solução para gestão territorial precisa considerar realidade de municípios pequenos
  \item \textbf{Múltiplas Personas}: Uma mesma ferramenta deve adaptar linguagem conforme o público
  \item \textbf{Integração com Realidade}: Ferramentas devem resolver problemas reais (não apenas teóricos)
  \item \textbf{Continuidade}: Projeto acadêmico precisa de planejamento para continuidade pós-defesa
\end{enumerate}

\newpage

% ==================== CONCLUSÕES ====================
\chapter{Conclusões}

\section{Objetivos Alcançados}

Este projeto de Iniciação Científica alcançou seus principais objetivos:

\begin{itemize}
  \item ✓ Desenvolver arquitetura RAG robusta para gestão territorial
  \item ✓ Implementar sistema web interativo com Google Sites
  \item ✓ Estruturar base de dados com 1000+ dados temáticos
  \item ✓ Validar solução com 9+ profissionais em 3 workshops
  \item ✓ Demonstrar viabilidade técnica e aceitação de usuários
  \item ✓ Documentar completamente o projeto
  \item ✓ Publicar em plataforma pública (Brasil Participativo)
\end{itemize}

\section{Contribuições Científicas}

\begin{enumerate}
  \item \textbf{Metodologia}: Abordagem de validação participativa para sistemas de IA em domínios especializados
  \item \textbf{Arquitetura}: Pattern de agentes especializados com orquestrador inteligente
  \item \textbf{Aplicação}: Demonstração de RAG em gestão territorial pública
  \item \textbf{Avaliação}: Metodologia de comparação RAG vs LLM puro com usuarios
\end{enumerate}

\section{Próximos Passos}

\subsection{Curto Prazo (Até Defesa de Mestrado - Nov/2025)}

\begin{itemize}
  \item Ativar API com credenciais UNIFESP
  \item Realizar último teste de sistema para banca
  \item Finalizar dissertação e artigo
  \item Preparar apresentação para banca
\end{itemize}

\subsection{Médio Prazo (Pós-Defesa - 2026)}

\begin{itemize}
  \item Publicar artigo em periódico científico
  \item Expandir base de dados com novos conteúdos
  \item Implementar sistema de feedback loop
  \item Otimizar para uso em produção
\end{itemize}

\subsection{Longo Prazo (2026+)}

\begin{itemize}
  \item Integração com sistemas governamentais
  \item Expansão para domínios relacionados (licenciamento, parcelamento)
  \item Modelos de computação visual
  \item Possível modelo de negócio/licenciamento
\end{itemize}

\section{Reflexão Final}

O projeto \textbf{Guia Interativo para Gestão Territorial com IA} representa uma ponte importante entre a pesquisa acadêmica e a solução de problemas práticos da administração pública. 

Demonstrou que:

\begin{quote}
  \textit{``Tecnologia de IA generativa, quando bem estruturada e validada com especialistas de domínio, pode ser uma ferramenta poderosa para democratizar acesso a conhecimento especializado e apoiar decisões mais informadas em gestão territorial.''}
\end{quote}

O projeto continuará evoluindo além desta Iniciação Científica, com perspectiva de impacto significativo para municípios e profissionais da área de gestão fundiária no Brasil.

\newpage

% ==================== REFERÊNCIAS ====================
\chapter*{Referências}

\begin{enumerate}
  \item LangChain Documentation: \url{https://python.langchain.com}
  \item LangGraph for Agent Orchestration: \url{https://github.com/langchain-ai/langgraph}
  \item FAISS (Facebook AI Similarity Search): \url{https://github.com/facebookresearch/faiss}
  \item FastAPI Documentation: \url{https://fastapi.tiangolo.com}
  \item Google AI Studio: \url{https://aistudio.google.com}
  \item Guia Interativo: \url{https://sites.google.com/view/guia-territorial/}
  \item Repositório GitHub: \url{https://github.com/Enzo0100/IC-UNIFESP-IA}
\end{enumerate}

\appendix

\chapter{Cronograma Executado}

\begin{table}[H]
\centering
\caption{Timeline Completo do Projeto}
\begin{tabular}{|l|l|l|}
\hline
\textbf{Período} & \textbf{Fase} & \textbf{Atividades Principais} \\
\hline
Maio/2025 & Planejamento & Requisitos, arquitetura, escolha de frameworks \\
\hline
Junho/2025 & MVP & Desenvolvimento backend, indexação, FAISS \\
\hline
Julho/2025 & Deploy & Deploy online, primeiros testes, ajustes iniciais \\
\hline
Agosto/2025 & Validação & 2 Workshops, feedback de usuários, otimizações \\
\hline
Setembro/2025 & Refinamento & Multi-agentes, feedback loop, último workshop \\
\hline
Outubro/2025 & Finalização & Documentação, submissão pública, preparação banca \\
\hline
\end{tabular}
\end{table}

\chapter{Arquivos Principais do Projeto}

\begin{verbatim}
IC-UNIFESP-IA/
├── README.md                          # Documentação principal
├── api.http                           # Exemplos de requisições HTTP
├── Relatorio_Final_IC.tex            # Este documento
├── conversas_whatsapp.txt            # Histórico de comunicação
├── backend/
│   ├── app.py                        # Aplicação principal FastAPI
│   ├── config.py                     # Configurações
│   ├── models.py                     # Modelos de dados
│   ├── api/
│   │   ├── endpoints.py             # Definição de endpoints
│   │   ├── lifespan.py              # Gerenciamento de ciclo de vida
│   │   └── __init__.py
│   └── utils/
│       ├── agents.py                # Definição de agentes
│       ├── cache.py                 # Cache de retrievers
│       ├── heuristics.py            # Heurísticas de decisão
│       ├── langgraph_pipeline.py    # Pipeline com LangGraph
│       ├── small_talk.py            # Detecção de small talk
│       ├── vectorstore.py           # Gerenciamento de FAISS
│       └── __init__.py
├── docker/
│   ├── Dockerfile                    # Imagem Docker
│   ├── docker-compose.yml           # Orquestração de containers
│   ├── requirements.txt              # Dependências Python
│   └── swarm.yml                     # Configuração Docker Swarm
├── frontend/
│   └── index.html                    # Interface HTML (descontinuada)
├── base_de_dados/
│   └── [Base de dados estruturada]   # Dados de treinamento
├── faiss_index/
│   └── index.faiss                   # Índice FAISS persistente
└── .gitignore, .env.example          # Configurações de projeto
\end{verbatim}

\end{document}
